\documentclass[10pt,a4paper]{amsart}
\usepackage[margin=19mm]{geometry}
\usepackage{amsmath,amssymb,mathtools}
\usepackage[scr=rsfs,cal=euler]{mathalfa}
\usepackage{xfrac}
\usepackage[unicode,hidelinks,bookmarks=false]{hyperref}
\newcommand{\ie}{i.e.\ }
\newcommand{\cf}{cf.\ }
\newcommand{\resp}{respectively\ }
\newcommand{\Subsection}{\subsection}
\providecommand{\nobreakdash}{\nobreak\hbox{-}}
\setlength{\emergencystretch}{2em}
\multlinegap0pt
\usepackage{amssymb}
\usepackage[all]{xy}
\usepackage{tikz}
\usepackage[normalem]{ulem}
\newtheorem{lem}{Lemma}
\newtheorem{thm}{Theorem}
\newtheorem{rem}{Remark}
\DeclareMathOperator{\FinSets}{FinSets}
\newcommand{\ca}[1]{\mathcal{#1}}
\DeclareMathOperator{\cat}{cat}
\title{Homotopy types of finite étale spaces and generalized inflations}
\author{Anton Ayzenberg, Nadya Khoroshavkina}
\date{Source: arXiv:2602.08694v1 (CC BY 4.0)}
\begin{document}
\maketitle

\noindent\emph{Source and licence.} Homotopy types of finite étale spaces and generalized inflations. Version arXiv:2602.08694v1 (CC BY 4.0), distributed by arXiv under the Creative Commons Attribution 4.0 International licence, http://creativecommons.org/licenses/by/4.0/, as recorded in the arXiv licence metadata for this exact version and shown on the abstract page https://arxiv.org/abs/2602.08694v1. Full licence notice: \texttt{licenses/arxiv-2602.08694-CC-BY-4.0.txt}.\par\smallskip

\noindent\textit{Editorial scope.} Counted unit: the article's lem lemInhFlabby (source0.tex lines 822--824). The article's complete original proof of it is delivered with it (source0.tex lines 826--835, one \texttt{proof} environment of 10 lines). No mathematics is written, repaired or re-derived: the statement and the proof are the article's own lines, in the article's own notation, and no retained line is substituted or re-flowed.  Retained uncounted as the source-local context the proof needs: lines 359--361; lines 702--704; lines 779--783.  Nothing was omitted from any retained span, so the package carries no omission record.  No retained line contains a citation, so the wrapper carries no bibliography and no bibliography entry.  No retained line contains a figure environment, an image-inclusion command or drawing code, so no image is delivered and the package declares no asset dependency.  Editorial formatting only, in addition: an AMS \texttt{amsart} preamble that restates the article's own declarations and macros used by the retained lines, namely the environments \texttt{lem}, \texttt{thm}, \texttt{rem} on separate counters, together with the standard packages those lines need.  The retained environments print in this document as Lemma 1, Theorem 1, Remark 1 in order of appearance, whereas the article prints the same items with the numbering of its own full document.  The complete native source of the article as distributed in the arXiv e-print for this exact version is bundled unchanged as \texttt{sources/194265/source0.tex}.
\begin{lem}[Gluing Lemma]\label{lemGlue}
Let $\ca{D}$ be a sheaf on a topological space $X_S$. Let $U_1,U_2$ be two open sets in $X_S$, such that for each point $x\in U_1\cap U_2$ the stalk $D(x)$ is a singleton. Then, for any two sections $v_1\in \ca{D}(U_1)$, $v_2\in \ca{D}(U_2)$, there exists a unique section $v\in \ca{D}(U_1\cup U_2)$ which restricts to $v_1$ on $U_1$ and to $v_2$ on $U_2$.
\end{lem}

\begin{thm}\label{thmSimplex}
Let $P=\ca{F}(\Delta^{n-1})$  and  $D\colon \cat(P^*)\to \FinSets$ be a diagram on the dual poset such that its corresponding sheaf $\ca{D}$ is inhabited  and  flabby. Then the geometric realization of the inflation $|P_D|$  is homotopy equivalent to a wedge of $(n-1)$-dimensional spheres.
\end{thm}

\begin{rem}\label{remOpenSubcomplexes}
Let, as in Theorem~\ref{thmSimplex}, $P=\ca{F}(\Delta^{n-1})$ be a poset of non-empty faces of a simplex  and  $D$ be a diagram on the poset $P^*$. The corresponding sheaf is given on the topological space $X_{P^*}$. Open sets in this topology are upper order ideals in $P^*$ or, equivalently, lower order ideals in $P$. A curious reader can notice that lower order ideals in the poset of faces of a simplex on a set of vertices $[n]$ are simplicial complexes on the same set $[n]$. An interesting interpretation emerges: open sets are simplicial subcomplexes. For a simplex $I\in P$ its minimal open neighbourhood $U_I$ is a simplicial subcomplex generated by this simplex, that is the one that contains $I$ and  all its faces.

Construction of a sheaf $\ca{D}$ on the topology $X_{P^*}$ is equivalent to assigning a finite set $\ca{D}(U)$ to each simplicial subcomplex $U$ in $\Delta^{n-1}$  and  a restriction morphism $\ca{D}(U\supset V)\colon \ca{D}(U)\to \ca{D}(V)$ for each pair of nested subcomplexes $U\supset V$.
\end{rem}

\begin{lem}\label{lemInhFlabby}
The sheaves $\ca{D}_1$ and $\ca{D}_2$, that correspond to diagrams $D_1$ and $D_2$, are inhabited and flabby.
\end{lem}

\begin{proof}
From the inhabitedness of $D$, it follows that all sets $D(J)$, for $J\in P$, are non-empty. The sets $A_1$, $A_2$ are non-empty by construction. All morphisms $D(J\geqslant I_0)$ are surjective, because they are special cases of the restriction morphisms in the sheaf $\ca{D}$, which is flabby. The inhabitedness of diagrams $D_1$, $D_2$ follows.

Let us prove the flabbiness of a sheaf $\ca{D}_\varepsilon$, $\varepsilon=1,2$. Let $U\supset V$ be two open sets in the Alexandrov topology $X_{P^*}$, i.e, according to Remark~\ref{remOpenSubcomplexes}, two simplicial subcomplexes. We need to prove that the map $\ca{D}_\varepsilon(U\supset V)\colon \ca{D}_\varepsilon(U)\to \ca{D}_\varepsilon(V)$ is surjective. Let us consider three cases.
\begin{enumerate}
  \item $I_0\notin U$. In this case the morphism $\ca{D}_\varepsilon(U\supset V)\colon \ca{D}_\varepsilon(U)\to \ca{D}_\varepsilon(V)$ literally coincides with the restriction morphism $\ca{D}(U\supset V)\colon \ca{D}(U)\to \ca{D}(V)$ of the original sheaf which is surjective by assumption.
  \item $I_0\in U$, $I_0\notin V$. Let $u\in \ca{D}_\varepsilon(V)=\ca{D}(V)$ be an arbitrary section on $V$. We consider an arbitrary section $v\in A_\varepsilon\subset D(I_0)=\ca{D}(U_{I_0})$. We recall that $U_{I_0}$ is a minimal open neighbourhood of $I_0$ in the Alexandrov topology and, at the same time, a simplicial subcomplex, generated by the simplex $I_0$, see Remark~\ref{remOpenSubcomplexes}. The choice of $I_0$ implies that, for each proper face $J<I_0$, the set $D(J)$ is a singleton. Therefore, the conditions of the Gluing Lemma~\ref{lemGlue} are satisfied for $U_{I_0}$ and $V$, which means there exists a section $u'\in \ca{D}(V\cup U_{I_0})$ on a subcomplex, generated by gluing a simplex $I_0$ to $V$. From the flabbiness of $\ca{D}$, it follows that there is a section $u''\in \ca{D}(U)$ extending $u'$ to the subcomplex $U$. We claim that $u''$ belongs to the subset of sections $\ca{D}_\varepsilon(U)\subset \ca{D}(U)$. This is practically a tautology: from the construction it follows that the restriction of $u''$ on the subcomplex $U_{I_0}$ is $v\in A_\varepsilon$, which actually means that the section $u''$ belongs to the specified sheaf $\ca{D}_\varepsilon$. So far, we proved that for any element $u\in \ca{D}_\varepsilon(V)$ there exists a covering element of it, namely $u''\in\ca{D}_\varepsilon(U)$, and therefore the map $\ca{D}_\varepsilon(U\supset V)$ is surjective.
  \item $I_0\in V\subset U$. This case can be proved analogously to the previous one, but the gluing procedure of $I_0$ to $V$ is unnecessary in this case.
\end{enumerate}
\end{proof}

\end{document}
